% Deckblatt
\thispagestyle{empty}
\vspace*{2cm}
\begin{center}
  \includegraphics[width=3cm]{images/TUMLogo_oZ_Outline_schwarz_CMYK}\\              \vspace{.5cm}
  {\Large \sc Technische Universität München}\\         \vspace{1cm}
  {\Huge \sc Fakultät für Informatik\\[1mm]}            \vspace{2cm}
  {\Large \textbf{Bachelorarbeit in Informatik}}\\      \vspace{1.5cm}
  {\huge \textbf{Intelligent Support}}\\ \vspace*{3mm}
  {\huge \textbf{for non-linear Serious Games}}\\               \vspace{1.5cm}
  Felix Kaser
  
  \vspace{1cm}
  \includegraphics[width=2.4cm]{images/IN_schwarz_CMYK}
\end{center}
%\rule{\textwidth}{1mm}
%\vspace{1cm}
%\noindent Version: \today, \hl{DRAFT}
\newpage
%%%%%%%%%%%%%%%%%%%%%%%%%%%%%%%
% Rückseite Deckblatt
\thispagestyle{empty}
\cleardoublepage
%%%%%%%%%%%%%%%%%%%%%%%%%%%%%%%
% Erste Seite (Titelblatt)
\thispagestyle{empty}
\begin{center}
  \includegraphics[width=3cm]{images/TUMLogo_oZ_Outline_schwarz_CMYK}\\              \vspace{.2cm}
  {\Large \sc Technische Universität München}\\                 \vspace{.2cm}
  {\huge \sc Fakultät für Informatik\\[1mm]}                    \vspace{1cm}
  {\Large \textbf{Bachelorarbeit in Informatik}}\\              \vspace{1cm}
  {\huge \textbf{Intelligent Support}}\\         \vspace*{3mm}
  {\huge \textbf{for non-linear Serious Games}}\\                       \vspace{1cm}
  {\huge \textbf{Intelligente Unterstützung}}\\   \vspace*{3mm}
  {\huge \textbf{nichtlinearer digitaler Lernspiele}}\\                     \vspace*{3mm}
  %~{\huge \textbf{in Verteilten Systemen}}\\                     \vspace{1cm}
    
  \parbox{1cm}{
  \begin{large}
  \begin{tabbing}
    Author: \hspace{2cm} \=Felix Kaser\\
    Supervisor: \>Prof. Bernd Brügge, PhD.\\
    Advisors: \>Dipl.-Inf. Univ. Dennis Pagano\\
    \>Dipl.-Inf. Univ. Damir Ismailović\\[1mm]
    Submission Date: \> March 15, 2010\\
  \end{tabbing}
  \end{large}
  }\\ %\vspace{.3cm}
  \includegraphics[width=2.4cm]{images/IN_schwarz_CMYK}
\end{center}
\newpage
\thispagestyle{empty}
\cleardoublepage
\thispagestyle{empty}
\begin{large}
\vspace*{13cm}
\noindent
I assure the single handed composition of this master's thesis only supported by declared resources.

\vspace{3cm}
\hspace*{3cm}%
Date, Signature: \dotfill
\\
\hspace*{9.5cm}%
\textit{(Felix Kaser)}
\end{large}
\newpage
\thispagestyle{empty}
\cleardoublepage
%abstract
\begin{center}
{\huge \textbf{Abstract}\\}
\end{center}
Learning games are the most promising trend in modern education. Educators try to
harvest the high motivation children show in computer games for learning. To achieve
this, many games have been developed: some are good, some are bad and most of them are something
in between. The winning strategy for such games is to respect the principles of pedagogics
and implement them in the game. Only a pedagogic sound game can effectively teach a matter.

This thesis analyzes the application domain and follows the steps of software
engineering to design a framework which -- in combination with a non-linear serious game
-- enhances the learning process. The framework monitors the actions of the learner
and detects if there are problems in the learning process which can be solved. This
is important for many different reasons: the learner does not get frustrated because
he is stuck somewhere in the game, playing the game makes learning more efficient,
the learner is more immersed in the game because he gets individual support by the game, etc.

At the end of this work, the reader will have a good overview of what state of
the art learning software is and how it can be improved using a framework which is designed to use
artificial intelligence in order to understand the behavior of the learner.

\newpage
\cleardoublepage
\begin{center}
{\huge \textbf{Dedication}\\}
\end{center}
\vspace{5cm}
\begin{quote}
\begin{center}
{\large To follow the path:\\
look to the master,\\
follow the master,\\
walk with the master,\\
see through the master,\\
become the master.\\[5mm]}
\hspace{4cm}
{\small \em A modern zen poem}
\end{center}
\end{quote}

I dedicate this work to all the ``masters'' of my life, who allowed me to learn from them.\\
%~My special thanks go to supervisors Dennis and Damir who helped me with their insight
%~which had a great impact on me. This work can never express how much I have learned
%~in the last 4 months.\\
%~I also want to thank my family and my friends, which supported me and helped me see the
%~light in the dark.
